\section*{अध्याय \ref{ch:b1:titration-methods} --- अनुमापन विधियाँ और वक्र}

\begin{solution}{exo:b1:titration-methods:1}
\ce{H3O+ + OH- -> 2H2O}, $K = 1/K_e = 10^{14.00}$। $C = 12.4 \times
0.100/20.0 = \qty{0.0620}{mol/L}$।
\end{solution}

\begin{solution}{exo:b1:titration-methods:2}
\omterm{def:b1:titration-methods:titration}{तुल्यता} पर विलयन \qty{0.050}{mol/L} पर अमोनियम क्लोराइड है:
$\mathrm{pH} = \frac12(9.25 - \log 0.050) = 5.28$।
मेथिल रेड (3.8--5.8) इसे समाहित करता है; मेथिल ऑरेंज, जिसका परास (2.5--4.5) pH गिरने पर \omterm{def:b1:titration-methods:titration}{तुल्यता} के बाद ही आता है, बहुत देर से रंग बदलेगा।
\end{solution}

\begin{solution}{exo:b1:titration-methods:3}
\ce{MnO4- + 5Fe^2+ + 8H+ -> Mn^2+ + 5Fe^3+ + 4H2O};
$C = 5 \times 0.0200 \times 12.6/10.0 = \qty{0.126}{mol/L}$।
\end{solution}

\begin{solution}{exo:b1:titration-methods:4}
$C = 0.0100 \times 14.2/50.0 = \qty{2.84e-3}{mol/L}$, कैल्सियम और
मैग्नीशियम मिलाकर। दोनों \omterm{def:b1:titration-methods:titration}{अंत्य बिंदु} से पहले अभिक्रिया करते हैं, जो केवल उनका योग देता है:
\omterm{def:b1:titration-methods:kinds}{युगपत् अनुमापन}।
\end{solution}

\begin{solution}{exo:b1:titration-methods:5}
0~mL: \omterm{def:b1:predominant-reaction:strong-acid}{दुर्बल अम्ल}, $\mathrm{pH} = \frac12(4.76 + 1.00) = 2.88$। 5.0~mL:
अर्ध-तुल्यता, 4.76। 10.0~mL: 8.73
(\cref{prop:b1:titration-methods:ph-at-equivalence})। 15.0~mL: आधिक्य
हाइड्रॉक्साइड, \qty{25.0}{mL} में \qty{0.50}{mmol}, $[\ce{OH-}] = 0.020$,
$\mathrm{pH} = 12.30$। सभी परिकलित वक्र से 0.01 तक मेल खाते हैं।
\end{solution}

\begin{solution}{exo:b1:titration-methods:6}
$\mathrm{p}K_a$ 2.15 और 7.21 में 4 से अधिक का अंतर है, और 7.21 तथा
12.34 में भी: दो \omterm{def:b1:titration-methods:titration}{तुल्यताएँ}, 10.0 और \qty{20.0}{mL} पर। पहली: \ce{H2PO4-} का
विलयन, $\mathrm{pH} \approx \frac12(2.15 + 7.21) = 4.68$; दूसरी:
\ce{HPO4^2-}, $\frac12(7.21 + 12.34) = 9.78$ (यथार्थ वक्र 4.71 और
9.66 देता है, सूत्र तनुकरण और पानी की उपेक्षा करते हैं)। तीसरी अम्लता
($\mathrm{p}K_a$ 12.34) पानी की अम्लता के बहुत पास है: मिलाया गया हाइड्रॉक्साइड
केवल आंशिक रूप से अभिक्रिया करता है और pH बिना किसी उछाल के बढ़ता है।
\end{solution}

\begin{solution}{exo:b1:titration-methods:7}
हाइड्रोक्लोरिक अम्ल की सांद्रता $6.2 \times 0.100/10.0 = \qty{0.062}{mol/L}$ है; एथेनोइक
अम्ल की $(14.6 - 6.2) \times 0.100/10.0 = \qty{0.084}{mol/L}$ है। ठीक बीच में
एथेनोइक अम्ल का आधा भाग अनुमापित होता है: $\mathrm{pH} = \mathrm{p}K_a = 4.76$।
\end{solution}

\begin{solution}{exo:b1:titration-methods:8}
$V_{\mathrm{eq}} = \qty{10.0}{mL}$। 2.0~mL: $E = 0.77 + 0.059\log(2/8) =
0.73$~V; 9.0~mL: $0.77 + 0.059\log 9 = 0.83$~V; 11.0~mL: आधिक्य
परमैंगनेट \qty{0.020}{mmol}, \ce{Mn^2+} \qty{0.200}{mmol}, $E = 1.51 +
\frac{0.059}{5}\log 0.10 = 1.50$~V; 20.0~mL: 1.51~V।
\end{solution}

\begin{solution}{exo:b1:titration-methods:9}
0~mL: $\kappa = (349.8 + 76.2) \times 0.0100 = \qty{4.26}{mS/cm}$।
10.0~mL: \ce{Na+} और \ce{Cl-}, $1.00/110 = \qty{9.09e-3}{mol/L}$ पर,
$\kappa = (50.3 + 76.2) \times \num{9.09e-3} = 1.15$। 20.0~mL: \ce{Na+}
$0.0167$, \ce{Cl-} और \ce{OH-} $0.00833$~mol/L, $\kappa = 0.84 + 0.635 + 1.65
= 3.13$~mS/cm। प्रवणताएँ (तनुकरण नगण्य, \qty{1}{mL}
\qty{1.0e-3}{mol/L} जोड़ता है): $(50.3 - 349.8) \times 10^{-3} = -0.30$ और
$(50.3 + 198.3) \times 10^{-3} = +0.25$~mS/cm प्रति mL।
\end{solution}

\begin{solution}{exo:b1:titration-methods:10}
$V_{\mathrm{eq}} - v$ पर आधिक्य अम्ल $C_Bv$ है, $V_{\mathrm{eq}} + v$ पर
आधिक्य क्षारक $C_Bv$; समान आयतन में पहले का $[\ce{H3O+}]$, बाद के
$[\ce{OH-}]$ के बराबर है, इसलिए $\mathrm{pH}(V_{\mathrm{eq}} - v) +
\mathrm{pH}(V_{\mathrm{eq}} + v) = 14$: (V$_{\mathrm{eq}}$,
7) के सापेक्ष सममिति। उछाल: \qty{20}{mL} में \qty{1.0e-5}{mol} आधिक्य, pH 3.30 से 10.70, 7.4
इकाई। 100 से भाग देने पर: \qty{20}{mL} में $\num{1.0e-7}$~mol, pH 5.3 से 8.7,
3.4 इकाई: कहीं कम तीक्ष्ण \omterm{def:b1:titration-methods:titration}{अंत्य बिंदु}।
\end{solution}

\begin{solution}{exo:b1:titration-methods:11}
पहले: बना \ce{Fe^3+} $= 5C_{\mathrm{Mn}}V$, बचा \ce{Fe^2+} $=
5C_{\mathrm{Mn}}(V_{\mathrm{eq}} - V)$, इसलिए $E = 0.77 + 0.059\log\frac{V}{V_{\mathrm{eq}}
- V}$: $V_{\mathrm{eq}}/2$ पर $E = 0.77$। बाद में: आधिक्य \ce{MnO4-} $\propto
V - V_{\mathrm{eq}}$, \ce{Mn^2+} $\propto V_{\mathrm{eq}}$: $E = 1.51 +
\frac{0.059}{5}\log\frac{V - V_{\mathrm{eq}}}{V_{\mathrm{eq}}}$ (pH 0);
$2V_{\mathrm{eq}}$ पर 1.51~V। \omterm{def:b1:titration-methods:titration}{तुल्यता} पर $(0.77 + 5 \times 1.51)/6 =
1.39$~V। pH 1 पर परमैंगनेट युग्म $0.059 \times 8/5 =
0.094$~V नीचे होता है: $E_{\mathrm{eq}} = (0.77 + 5 \times 1.416)/6 = 1.31$~V।
\end{solution}

\begin{solution}{exo:b1:titration-methods:12}
\ce{NH4+ + OH- -> NH3 + H2O} का $K = 10^{14.00 - 9.25} = 10^{4.75}$ है:
\omterm{def:b1:titration-methods:titration}{तुल्यता} के पास अभिक्रिया पूर्ण नहीं होती और उछाल छोटा होता है,
pH 11 के आसपास (\omterm{def:b1:titration-methods:titration}{तुल्यता} pH $14 - \frac12(4.75 + 1.30) = 11.0$)। \omterm{def:b1:titration-methods:kinds}{पश्च
अनुमापन}: सोडियम हाइड्रॉक्साइड की $n_1$ मात्रा (आधिक्य) मिलाइए, बनी अमोनिया को
उबालकर निकाल दीजिए, फिर बचे हाइड्रॉक्साइड का हाइड्रोक्लोरिक अम्ल से \omterm{def:b1:titration-methods:titration}{अनुमापन} कीजिए
($n_2$, तीक्ष्ण प्रबल--प्रबल उछाल): $n(\ce{NH4+}) = n_1 - n_2$।
\end{solution}

\begin{solution}{pb:b1:titration-methods:1}
\textbf{1.} \ce{C6H6O6 + 2H+ + 2e- -> C6H8O6}।
\textbf{2.} \ce{C6H8O6 + I2 -> C6H6O6 + 2I- + 2H+}।
\textbf{3.} 0 और $-$I।
\textbf{4.} \ce{I2 + 2S2O3^2- -> 2I- + S4O6^2-}, $\log K = 2(0.53 - 0.02)/0.059
= 17$।
\textbf{5.} नीले रंग के लिए आयोडीन चाहिए; यह अंतिम
आयोडीन के साथ, \omterm{def:b1:titration-methods:titration}{तुल्यता} पर, लुप्त हो जाता है।
\textbf{6.} वायु से हानि परिणाम को घटा देती; तुरंत और
आधिक्य में मिलाई गई आयोडीन ऐस्कॉर्बिक अम्ल को जल्दी और पूरी तरह ऑक्सीकृत कर देती है।
\textbf{7.} पूरे \omterm{def:b1:titration-methods:titration}{अनुमापन} के दौरान एक तीव्र, \omterm{def:b1:extent-q-and-k:quantitative}{मात्रात्मक अभिक्रिया},
\omterm{def:b1:titration-methods:titration}{अनुमापक} को धीरे-धीरे मिलाते समय वायु द्वारा ऑक्सीकरण के बिना, और आयोडीन के पहले
आधिक्य पर एक \omterm{def:b1:titration-methods:titration}{अंत्य बिंदु}।
\textbf{8.} आयोडीन का आधिक्य $n_R$; ऐस्कॉर्बिक अम्ल से अभिक्रिया; बची हुई
आयोडीन का थायोसल्फ़ेट से \omterm{def:b1:titration-methods:titration}{अनुमापन}।
\textbf{9.} अन्यथा कुछ ऐस्कॉर्बिक अम्ल बचा रह जाता और आयोडीन का
आधिक्य, शून्य, यह नहीं बताता कि कितना। यहाँ
\qty{5.000e-4}{mol} में से \qty{2.775e-4}{mol} ने अभिक्रिया की: सचमुच आधिक्य।
\textbf{10.} थायोसल्फ़ेट विलयन समय के साथ धीरे-धीरे बदलते हैं; इसकी
सांद्रता सीधे परिणाम में प्रवेश करती है।
\textbf{11.} 10।
\textbf{12.} \qty{20.00}{mL} का आयतनमितीय पिपेट।
\textbf{13.} $20.00 \times 10^{-3} \times 0.0250 = \qty{5.000e-4}{mol}$।
\textbf{14.} $8.90 \times 10^{-3} \times 0.0500 = \qty{4.450e-4}{mol}$।
\textbf{15.} आधा: \qty{2.225e-4}{mol}।
\textbf{16.} $5.000 - 2.225 = \qty{2.775e-4}{mol}$।
\textbf{17.} एक-एक के अनुपात में: \qty{2.775e-4}{mol}।
\textbf{18.} $\times 10$: \qty{2.775e-3}{mol}, $\times 176.0$:
\qty{0.488}{g}।
\textbf{19.} लेबल से लगभग \qty{2}{\percent} कम।
\textbf{20.} $n = \bigl(C_IV_I - \frac12C_TV_T\bigr)\,V_{\mathrm{flask}}/V_{\mathrm{aliquot}}$।
\textbf{21.} मिलाई गई आयोडीन: पिपेट $0.03/20.00 = \qty{0.15}{\percent}$
और सांद्रता \qty{0.2}{\percent}, लगभग \qty{0.25}{\percent}, अर्थात्
\qty{1.3e-6}{mol}। आधिक्य: ब्यूरेट $0.05/8.90 = \qty{0.56}{\percent}$ और
सांद्रता \qty{0.2}{\percent}, लगभग \qty{0.60}{\percent}, अर्थात्
\qty{1.3e-6}{mol}।
\textbf{22.} निरपेक्ष अनिश्चितताएँ जुड़ती हैं, जबकि अंतर
दोनों में से किसी भी मात्रा से छोटा है: $\sqrt{1.3^2 + 1.3^2} \times 10^{-6} =
\qty{1.8e-6}{mol}$, \qty{2.775e-4}{mol} का \qty{0.66}{\percent}, जो
किसी भी एक पद से अधिक है।
\textbf{23.} फ़्लास्क (\qty{0.1}{\percent}) और भाग वाले पिपेट
(\qty{0.2}{\percent}) के साथ: $\sqrt{0.66^2 + 0.1^2 + 0.2^2} = \qty{0.70}{\percent}$।
\textbf{24.} $0.488 \pm 0.003$~g: \textbf{प्रति गोली \qty{0.49}{g} ऐस्कॉर्बिक
अम्ल}।
\end{solution}
